% Part: first-order-logic
% Chapter: syntax-and-semantics
% Section: free-vars-sentences

\documentclass[../../../include/open-logic-section]{subfiles}

\begin{document}

% SEGMENT OLP-0157-S01
\olfileid{fol}{syn}{fvs}

\olsection{கட்டற்ற \printtoken{P}{variable} மற்றும் \printtoken{P}{sentence}}

\begin{defn}[!!a{variable} இன் கட்டற்ற நிகழ்வுகள்]
\ollabel{defn:free-occ}
ஒரு !!a{formula} இல் !!a{variable} இன் \emph{கட்டற்ற} நிகழ்வுகள்
பின்வருமாறு தொகுத்தறிதல் வழியில் வரையறுக்கப்படுகின்றன:
\begin{enumerate}
\item \indcase*{!A}{$!A$ is atomic}{$\indfrm$ இல் உள்ள அனைத்து !!{variable}
  நிகழ்வுகளும் கட்டற்றவை.}

\tagitem{prvNot}{\indcase{!A}{\lnot !B}{$\indfrm$ இன் கட்டற்ற !!{variable}
  நிகழ்வுகள், $!B$ இல் உள்ளவற்றுடன் சரியாக ஒத்தவை.}}{}

\item \indcase{!A}{(!B \ast !C)}{$\indfrm$ இன் கட்டற்ற !!{variable}
  நிகழ்வுகள் $!B$ இல் உள்ளவையும்~$!C$ இல் உள்ளவையும் ஆகும்.}

\tagitem{prvAll}{\indcase{!A}{\lforall[x][!B]}{$\indfrm$ இல் உள்ள கட்டற்ற
  !!{variable} நிகழ்வுகள், $!B$ இல் உள்ள $x$ நிகழ்வுகளைத் தவிர மற்றவை.}}{}

\tagitem{prvEx}{\indcase{!A}{\lexists[x][!B]}{$\indfrm$ இல் உள்ள கட்டற்ற
  !!{variable} நிகழ்வுகள், $!B$ இல் உள்ள $x$ நிகழ்வுகளைத் தவிர மற்றவை.}}{}
\end{enumerate}
\end{defn}

\begin{defn}[கட்டுண்ட மாறிகள்]
ஒரு சூத்திரம்~$!A$ இல் !!a{variable} இன் ஒரு நிகழ்வு \emph{கட்டுண்டது}
எனில் அது கட்டற்றதல்ல.
\end{defn}

\begin{prob}
\olref[fol][syn][fvs]{defn:free-occ} இன் வழியில் கட்டுண்ட மாறி
நிகழ்வுகளுக்கான ஒரு தொகுத்தறிதல் வரையறையைத் தருக.
\end{prob}

\begin{defn}[வீச்சு]
\iftag{prvAll}{ஒரு சூத்திரம்~$!A$ இல் உள்ள துணைச்சூத்திரத்தின் நிகழ்வு
  $\lforall[x][!B]$ ஆக இருந்தால், $!A$ இல் அதற்குரிய $!B$ நிகழ்வு,
  $\lforall[x]$ இன் அதற்குரிய நிகழ்வின் \emph{வீச்சு}. 
  \iftag{prvEx}{$\lexists[x]$ க்கும் இதேபோல்.}{}}{ஒரு
  சூத்திரம்~$!A$ இல் $\lexists[x][!B]$ துணைச்சூத்திரத்தின் நிகழ்வு
  இருந்தால், $!A$ இல் அதற்குரிய $!B$ நிகழ்வு, $\lexists[x]$ இன்
  அதற்குரிய நிகழ்வின் \emph{வீச்சு}.}

$!B$ ஒரு அளவையடை நிகழ்வின் வீச்சாக இருந்தால்
\iftag{prvAll}{$\lforall[x]$\iftag{prvEx}{ அல்லது
    $\lexists[x]$}{}}{$\lexists[x]$} இல்~$!A$, அப்போது~$!B$ இல் $x$ இன்
கட்டற்ற நிகழ்வுகள்
\iftag{prvAll}{$\lforall[x][!B]$\iftag{prvEx}{ மற்றும்
$\lexists[x][!B]$}{}}{$\lexists[x][!B]$} இல் கட்டுண்டவை. இந்நிகழ்வுகள்
குறிப்பிட்ட அளவையடை நிகழ்வால் \emph{கட்டுண்டவை} என்று சொல்கிறோம்.
\end{defn}

\begin{ex}
பின்வரும் சூத்திரத்தைக் கருதுக:
\[
\lexists[\Obj v_0][\underbrace{\Atom{\Obj A^2_0}{\Obj v_0,\Obj v_1}}_{!B}] 
\]
$!B$ என்பது $\lexists[\Obj v_0]$ இன் வீச்சு. இந்த அளவையடை $!B$ இல்
$\Obj v_0$ இன் நிகழ்வைக் கட்டுகிறது; ஆனால் $\Obj v_1$ இன் நிகழ்வைக்
கட்டுவதில்லை. எனவே இந்நிலையில் $\Obj v_1$ ஒரு கட்டற்ற மாறி.


இது இன்னும் சிக்கலான !!{formula}~$!A$ இல் எவ்வாறு இயங்கும் என்பதைப்
பார்க்கலாம்:
\[
\lforall[\Obj v_0][\underbrace{(\Atom{\Obj A^1_0}{\Obj v_0} \lif
    \Atom{\Obj A^2_0}{\Obj v_0, \Obj v_1})}_{!B}] \lif \lexists[\Obj
  v_1][\underbrace{(\Atom{\Obj A^2_1}{\Obj v_0, \Obj v_1} \lor \lforall[\Obj v_0][\overbrace{\lnot \Atom{\Obj A^1_1}{\Obj v_0}}^{!D}])}_{!C}]
\]
$!B$ என்பது முதல் $\lforall[\Obj v_0]$ இன் வீச்சு; $!C$ என்பது
$\lexists[\Obj v_1]$ இன் வீச்சு; $!D$ என்பது இரண்டாவது
$\lforall[\Obj v_0]$ இன் வீச்சு. முதல் $\lforall[\Obj v_0]$,
$!B$ இல் உள்ள $\Obj v_0$ நிகழ்வுகளைக் கட்டுகிறது; $\lexists[\Obj v_1]$,
$!C$ இல் உள்ள $\Obj v_1$ நிகழ்வைக் கட்டுகிறது; இரண்டாவது
$\lforall[\Obj v_0]$, $!D$ இல் உள்ள $\Obj v_0$ நிகழ்வைக் கட்டுகிறது.
$\Obj v_1$ இன் முதல் நிகழ்வும் $\Obj v_0$ இன் நான்காவது நிகழ்வும் $!A$ இல்
கட்டற்றவை. $\Obj v_0$ இன் கடைசி நிகழ்வு $!D$ இல் கட்டற்றது; ஆனால் $!C$
மற்றும்~$!A$ இல் கட்டுண்டது.
\end{ex}

\begin{defn}[வாக்கியம்]
!!^a{formula}~$!A$ ஒரு \article{sentence} \emph{!!{sentence}} ஆகும்
அப்போதும் அப்போது மட்டுமே, அதில் !!{variable}s இன் கட்டற்ற நிகழ்வுகள்
எதுவும் இல்லாவிட்டால்.
\end{defn}

% add examples!

\end{document}
